\documentclass[12pt]{article}
\usepackage[margin=0.58in,top=0.62in,bottom=0.48in]{geometry}
\usepackage{lmodern}
\usepackage{microtype}
\usepackage{fancyhdr}
\usepackage{titlesec}
\usepackage{enumitem}
\usepackage{lastpage}
\usepackage[hidelinks]{hyperref}

\setlength{\parindent}{0pt}
\setlength{\parskip}{0.18em}
\setlength{\headheight}{26pt}
\raggedbottom
\titleformat{\section}{\normalsize\bfseries\scshape}{}{0pt}{}
\titleformat{\subsection}{\normalsize\bfseries}{}{0pt}{}
\titlespacing*{\section}{0pt}{0.35em}{0.12em}
\titlespacing*{\subsection}{0pt}{0.25em}{0.08em}
\pagestyle{fancy}
\fancyhf{}
\fancyhead[C]{\footnotesize Lit Example Reader \quad Assignment ID: U1L12LE \quad Page \thepage\ of \pageref{LastPage}}
\renewcommand{\headrulewidth}{0.5pt}

\newenvironment{playpassage}{\begin{quote}\footnotesize\setlength{\parskip}{0.08em}\setlength{\parindent}{0pt}}{\end{quote}}
\newcommand{\speaker}[1]{\textbf{#1.}}

\begin{document}

\begin{center}
{\LARGE\bfseries Week 12 Lit Example Reader}\par\vspace{0.02em}
{\large\scshape Shakespeare: Examples, Character Claims, and Overclaims}\par\vspace{0.06em}
\rule{0.78\textwidth}{0.5pt}
\end{center}

\textbf{Purpose:} Shakespeare passages for Week 12: test whether observed examples justify a general claim about a person.

\textbf{What to watch for:} Examples, pattern, sample size, missing cases, counterexamples, and revised claims.

\section*{Before the Passages: The Week 12 Logic Tool}

\begin{itemize}[leftmargin=1.3em,itemsep=0.05em,topsep=0.08em,parsep=0.03em]
\item \textbf{Example:} one observed case or piece of evidence.
\item \textbf{General claim:} a broader conclusion about character, motive, or pattern.
\item \textbf{Overclaim:} a conclusion larger than the evidence supports.
\item \textbf{Revised claim:} a narrower conclusion that fits the evidence.
\end{itemize}

\section*{Passage 1: Julius Caesar, Act 1, Scene 2 --- Cassius Builds a Case from Anecdotes}

\textbf{Context:} Cassius wants Brutus to see Caesar as weak and unworthy of Rome's awe. Watch how a few personal examples become a broad character claim.

\begin{playpassage}
\speaker{Cassius} I was born free as Caesar; so were you:\\
We both have fed as well, and we can both\\
Endure the winter's cold as well as he:\\
For once, upon a raw and gusty day,\\
The troubled Tiber chafing with her shores,\\
Caesar said to me, ``Darest thou, Cassius, now\\
Leap in with me into this angry flood,\\
And swim to yonder point?'' Upon the word,\\
Accoutred as I was, I plunged in\\
And bade him follow; so indeed he did.\\
...\\
Caesar cried, ``Help me, Cassius, or I sink!''\\
...\\
He had a fever when he was in Spain,\\
And when the fit was on him, I did mark\\
How he did shake: 'tis true, this god did shake.
\end{playpassage}

\subsection*{Reader Notes}

This is your assisted example. Cassius offers particular cases: Caesar needed help swimming, and Caesar shook during a fever. These examples may support a limited claim: Caesar is physically vulnerable and not literally godlike. They do not by themselves prove every broader claim Cassius wants Brutus to feel, such as ``Caesar is unfit to lead'' or ``Rome must resist him.'' Week 12 asks how far the examples actually reach.

\textbf{Reading task:} Number Cassius's two examples. In the margin, write one limited claim they support and one larger claim they do not yet prove.

\newpage

\section*{Passage 2: Macbeth, Act 1, Scene 5 --- Lady Macbeth Reads Macbeth's Character}

\textbf{Context:} Lady Macbeth has read Macbeth's letter about the witches and the new title. She now makes a broad judgment about Macbeth's ambition and nature.

\begin{playpassage}
\speaker{Lady Macbeth} Glamis thou art, and Cawdor; and shalt be\\
What thou art promised: yet do I fear thy nature;\\
It is too full o' the milk of human kindness\\
To catch the nearest way: thou wouldst be great;\\
Art not without ambition, but without\\
The illness should attend it.
\end{playpassage}

\subsection*{Reader Notes}

This passage shifts more work to you. Lady Macbeth is not using a public survey or a scientific sample. She knows Macbeth intimately, but the passage gives the reader only a small window: the letter, the prophecy, and her judgment of his nature. Mark what evidence is visible in the passage. Then ask what claim about Macbeth it supports, and what stronger claim would require more evidence from the play.

\textbf{Reading task:} Underline Lady Macbeth's general claim about Macbeth. Next to it, write whether the visible evidence is one example, several examples, or a long pattern. Save the final revision for the worksheet.

\section*{How to Use These Passages on the Worksheet}

Do not summarize the plot. Use Week 12's evidence ladder: example, pattern, limited claim, overclaim, revised claim.

\end{document}
